\chapter{Conclusiones y Trabajo Futuro}
\label{ch:conclusiones}

\section{Síntesis de hallazgos}
\label{sec:concl-sintesis}

Esta tesis se propuso romper la circularidad que socava la planeación de transporte
basada en datos y construir, en su lugar, una metodología basada en la demanda que
diagnostica la necesidad no atendida, genera corredores candidatos y se transfiere a
otras ciudades. Las cuatro preguntas de investigación pueden responderse ahora de forma
directa.

\paragraph{PI1 -- Demanda.} \emph{Sí.} La demanda de transporte puede estimarse a
resolución de \ageb solo a partir de datos de Nivel~1. La generación de viajes y el modelo
gravitacional doblemente restringido del \cref{sec:w1} producen una superficie de demanda
metropolitana usando únicamente variables censales, económicas y de red, sin ningún insumo
de transporte, y la calibración contra la encuesta EOD~2022 (\cref{sec:w2}) fija el
parámetro de decaimiento en $\betadecay=1.2005$ ---un valor que ajusta los flujos
observados mejor que el previo convencional. La demanda, por tanto, nunca es función de la
oferta actual, que es la precondición de la que depende el resto del marco.

\paragraph{PI2 -- Diagnóstico.} \emph{Sí, y se valida fuera de muestra.} El índice de
brecha de cobertura (\cref{sec:w3}) contrapone la demanda modelada con la accesibilidad
medida por \textsc{gtfs}, produciendo una variable dependiente construida con independencia
de la oferta; el reentrenamiento supervisado la aprende de las características de lugar sin
fuga (\cref{tab:retrain}). Su mayor respaldo es el experimento natural de la Línea~4: un
corredor que el modelo nunca vio, luego elegido para construcción, muestra \num{1.6} veces
la tasa metropolitana de brecha alta. Este diagnóstico es la contribución más firme y
citable de la tesis.

\paragraph{PI3 -- Generación.} \emph{Parcialmente, y ahora caracterizada.} Una vez que los
corredores se modelan como trayectorias reales (un tronco por diámetro del árbol de
expansión mínima sobre el grafo vial) y se juzgan por rectitud respecto a las anclas en vez
de por el rodeo entre extremos (\cref{sec:w6}), el generador produce al menos un corredor
sustantivo, factible y que aprueba en mérito ---\code{W6\_G02}, \SI{12.1}{\kilo\metre},
\SI{56}{\percent} de brecha alta, no redundante, eficiencia de demanda en el percentil
\num{73}. También arroja conectores de menor eficiencia junto a él, y no reconstruye las
líneas férreas ya construidas ---un resultado esperado, pues la reconstrucción es un proxy
débil del mérito (\cref{sec:disc-reconstruction}). La capa generativa es así una
contribución genuina pero acotada.

\paragraph{PI4 -- Transferibilidad.} \emph{Sí, de extremo a extremo.} La metodología
completa W1--W8 corre sin rediseño a la medida en dos zonas metropolitanas contrastantes
---Toluca (grande, dieciséis municipios) y Aguascalientes (compacta, tres)---
re-derivando demanda, diagnóstico, pesos, corredores y auditoría a partir de los datos de
Nivel~1 y \textsc{gtfs} propios de cada ciudad (\cref{ch:transferibilidad}). El diagnóstico
diferencia sensatamente a las tres metrópolis por su proporción de demanda no atendida, y
la ponderación de cada ciudad se re-deriva de sus propios datos en lugar de heredar la de
la \zmg.

\section{Recomendaciones prácticas}
\label{sec:concl-recomendaciones}

Para una agencia de planeación, el marco es más útil como \emph{tamiz diagnóstico} que como
piloto automático. De ahí se siguen tres recomendaciones. Primera, tratar la superficie de
brecha de cobertura como el producto principal: ordena dónde se concentra la demanda no
atendida a una resolución más fina que cualquier línea individual, y puede recalcularse a
bajo costo conforme se actualizan los datos censales y \textsc{gtfs}. Segunda, usar la
auditoría de rutas (\cref{sec:w7}) para marcar rutas existentes de baja demanda, indirectas
y redundantes para su revisión ---en las ciudades de transferencia sacó a la luz una
redundancia sustancial de servicio paralelo, candidata a racionalización. Tercera, leer los
corredores generados como \emph{hipótesis a estudiar}, no como trazos finales: el corredor
que aprueba en mérito es un punto de partida defendible para un estudio de factibilidad,
pero las limitaciones caracterizadas del generador implican que su salida debe informar, no
reemplazar, el proceso de ingeniería y comunidad.

\section{Trabajo futuro}
\label{sec:concl-futuro}

Cuatro direcciones fortalecerían más el marco. Primera, la \emph{calibración local para las
ciudades de transferencia}: Toluca y Aguascalientes heredan hoy el previo de decaimiento de
la \zmg, y un $\betadecay$ calibrado con encuesta por ciudad afinaría sus superficies de
demanda. Segunda, un \emph{modelo gravitacional con distancia de red}: reemplazar la
impedancia euclidiana por impedancia de tiempo de viaje sobre el grafo vial mejoraría las
magnitudes absolutas de demanda. Tercera, \emph{transferir la capa supervisada}: el
reentrenamiento de brecha de cobertura se corrió solo para la \zmg, y aplicarlo a las
ciudades de transferencia probaría si la relación aprendida entre características de lugar y
demanda no atendida es en sí misma portable. Cuarta, un \emph{modelador híbrido de
corredores} que alterne entre el tronco por diámetro y el ruteo tipo agente viajero donde
este último siga siendo factible, lo que el análisis exploratorio sugiere que mejoraría
marginalmente la cobertura sin sacrificar factibilidad. Más allá de estas, la validación
externa más clara vendría de la afluencia observada en cualquier corredor que llegue a
construirse a lo largo de un trazo generado ---el contrafactual que, por construcción,
ningún estudio puramente ex ante puede aportar.
